\chapter{Reconstruction Method}

Classical light-field rendering combines and resamples captured views to synthesize a view from a new viewpoint \cite{levoyHanrahanLightFieldRendering2023}. This thesis takes the view-combination idea as a starting point, while interpreting the captured sub-aperture images according to the Fourier light-field microscopy (FLFM) image-formation model of Guo et al. \cite{guo2019fourierLightField}. The goal here is different from novel-view synthesis: to reconstruct a single two-dimensional image from the views of an FLFM acquisition. To this end, we propose a two-stage process that first aligns and combines the views globally, then refines them locally pixel-wise before combining them again.

To mitigate optical aberrations associated with scattering media, the reconstruction uses two complementary stages. Both use the shared scene information in the sub-aperture images to estimate shifts before averaging, but they differ in the spatial scale of the correction. The global stage assigns one uniform "global" translation to each view, bringing the views into broad alignment and producing an initial average. This corrects overall offsets, but cannot account for aberrations whose displacement varies across a single view. The local stage therefore follows by estimating a separate shift at each image location in each view, using the globally aligned views as its starting point. This finer correction can address residual, spatially varying shifts before the views are averaged again. Applying the global correction first establishes broad correspondence; the local correction then refines it. The stages thus share a matching-and-averaging strategy, while addressing different scales of image displacement.

In the global stage, the central sub-aperture image (view) serves as the reference. To estimate each view's translation, the method first creates thresholded copies of the reference and the moving view: pixels below a chosen fraction of their own maximum intensity are set to zero. This step is intended to reduce the impact of noise, defocused light, and highly-distorted features. The raw cross-correlation of the copies is computed, over a bounded range of integer shifts, and the shift with the largest correlation is assigned to the moving view.

The estimated translations are applied to the original, unthresholded sub-aperture images; thresholding is used only to determine the shifts. The output of the global stage is the arithmetic mean of the aligned views, including the unchanged reference view. Thus, all views contribute equally, regardless of their correlation score or the amount of overlap after shifting. This provides an initial reconstruction (referred to as the global average), in which the views are aligned as wholes, while leaving spatially varying displacement for the local stage.

In the local stage, each globally aligned view is compared with the global average, which serves as the new reference. For each view, the method tests a bounded set of possible shifts at every specific pixel and selects the shift with the strongest local match. Matching is performed on thresholded copies of the view and reference; the original view intensities are retained for the reconstruction. When the images are upsampled for matching, the tested shifts can be finer than one pixel.

The local match is based on the product of the view and reference intensities, averaged across the local neighborhood with greater weight near the pixel being evaluated and normalized by their local intensities. This gives a correlation-like score, rather than the raw cross-correlation used in the global stage. Gaussian smoothing lets nearby pixels contribute evidence to the choice of shift, motivated by the expectation that aberration varies gradually across the image, so nearby pixels may favor similar shifts. The smoothing does not force them to choose the same shift; the best shift is selected independently at each location. A separate option can smooth the resulting shift maps, but it is disabled by default.

The selected shifts are used to resample each globally aligned view: for every output pixel, the method reads the corresponding location in the aligned view, using linear interpolation. Samples beyond the view boundary are set to zero. The resampled views are then averaged to form the local stage output, referred to as the local average. In this way, the local stage refines alignment within each view before combining the views, while retaining the same overall matching-and-averaging principle as the global stage.
